\section{Fourier Analysis}

	The \textit{Fourier transform} $\mathcal{F}$ is an operator (an isometric isomorphism) mapping one function to another. There are three spaces where it is naturally defined, the \textit{Schwartz class} $\mathcal{S}$, the $L^2$ space, and the Banach dual of $\mathcal{S}$. 
	\begin{align*}
			\mathcal{F} &: \mathcal{S} \longrightarrow \mathcal{S} \\
			\mathcal{F} &: L^2 \longrightarrow L^2 \\
			\mathcal{F} &: \mathcal{S}^\prime \longrightarrow \mathcal{S}^\prime 
	\end{align*}
	and $\mathcal{S} \subset L^p \subset \mathcal{S}^\prime$. $\mathcal{S}$ is small, $\mathcal{S}^\prime$ is big, and $L^p$ spaces live in between them. The \textit{Schwartz class} is the function space of all functions $f: \mathbb{R}^n \longrightarrow \mathbb{C}$ satisfying 
	\begin{enumerate}
			\item $f \in C^\infty (\mathbb{R}^n, \mathbb{C})$ 
			\item $f$ decays rapidly in the sense that for any scalar $\alpha$, $f$ multiplied by any monomial $x^\alpha$ is bounded everywhere: 
			\[\sup_{x \in \mathbb{R}^n} | x^\alpha f (x)| < \infty \text{, or a similar condition } ||x^\alpha f ||_\infty < \infty \]
			Furthermore, with vector $\beta = (\beta_1, \ldots, \beta_n) \in \mathbb{N}^n$ and multi-indexing $D^\beta = D_1^{\beta_1} D_2^{\beta_2} \ldots D_n^{\beta_n} f$, all of its derivatives are also bounded for all $\beta$
			\[\sup_{x \in \mathbb{R}^n} |x^\alpha D^\beta f(x)| < \infty, \text{ or } ||x^\alpha D^\beta f| < \infty\]
	\end{enumerate}
	Now Schwartz functions have both the properties of smoothness and decay. The purpose of the Fourier transform is to convert between the properties of smoothness and decay. That is, if $f$ decays (is smooth) to some extent, then $\mathcal{F}f$ is smooth (decays) to that extent. $\mathcal{F}: \mathcal{S} \longrightarrow \mathcal{S}$ is defined (note that $x$ and $\xi$ are just dummy variables)
	\[f(x) \mapsto \mathcal{F} f (\xi) = \hat{f} (\xi) = \int_{\mathbb{R}^n} f(x) \, e^{-i x \xi} \, dx \]
	This transform is a simple map from a larger class called the \textit{Fourier integral operators}, which are maps of the form 
	\[f(x) \mapsto \mathcal{F} f (\xi) = \int_{\mathbb{R}^n} f(x) \, e^{\varphi(x, \xi)} a(\xi, x)\, dx\]
	Now since $\mathcal{F}$ is an isometric automorphism of $\mathcal{S}$, we can take its inverse transform, defined 
	\[\mathcal{F}^{-1} g(x) = \int_{\mathbb{R}^n} g(\xi) \, e^{i x \xi} \, d\xi\]
	The Fourier transform allows us to connect multiplication and differentiation. Let us calculate the $\alpha$-fold derivative along $\xi$ of the Fourier transform of $f$. 
	\begin{align*}
			\partial_\xi^\alpha \hat{f}(\xi) & = \partial_\xi^\alpha \int_{\mathbb{R}^n} e^{-i x \xi} f(x)\,dx \\
			& = \int_{\mathbb{R}^n} \partial_\xi^\alpha e^{-i x \xi} f(x)\,dx \\
			& = \int_{\mathbb{R}^n} (-i x)^\alpha e^{-i x \xi} f(x) \\
			& = \mathcal{F}((-ix)^\alpha f(x)) 
	\end{align*}
	So, a $\alpha$-fold differentiation after Fourier transforming a function is the same as multiplying $x^\alpha$ first and then Fourier transforming the same function. Now, if we multiply a the Fourier transform of a function by a monomial this sort of looks like the Fourier transform of the derivative of the function. 
	\begin{align*}
			\xi^\alpha \hat{f}(\xi) & = \int_{\mathbb{R}^n} \xi^\alpha e^{-i x \xi} \, f(x) \, dx \\
			& = \int_{\mathbb{R}^n} \partial_x^\alpha (e^{-i x \xi}) \, f(x)\, dx \\
			& = \int_{\mathbb{R}^n} (e^{-i x \xi}) \, \partial_x^\alpha f(x)\, dx & (\text{Integration by Parts})
	\end{align*}

	Minkowski Inequality says that if you have an $L^p(\Omega; \mathbb{R})$ space with norm $|| \cdot ||_{L^p} \coloneqq ( \int |f|^p )^{1/p}$, then the norm satisfies the triangle inequality. Using induction gives
	\[\bigg| \bigg| \sum_j f_j (x) \bigg| \bigg|_{L^p (dx)} \leq \sum_j \big| \big| f_j (x) \big| \big|_{L^p (dx)}\]
	That is, the norm of the sum of the functions is less than the sum of the function norms. But if we index $f_j (x)$ as $f(x, y)$, then since integration is really the same thing as summation, we can write the above as 
	\[\bigg| \bigg| \int_{\mathbb{N}} f(x, y) \,dy \bigg| \bigg|_{L^p (dx)} \leq \int_\mathbb{N} \big| \big| f(x, y) \big| \big|_{L^p (dx)} \,dy \]
	expanding 
	\[\bigg( \int_\Omega \bigg| \int_{\mathbb{N}} f(x, y) \,dy \bigg|^p \, dx \bigg)^{\frac{1}{p}} \leq \int_\mathbb{N} \bigg( \int_\Omega \big| f(x, y) \big|^{p} \,dx \bigg)^{\frac{1}{p}} \, dy\]
	So summing $\int\,dy$, then power to $p$, then norming $\int\,dx$, then power to $\frac{1}{p}$ is less than powering to $p$ first, then norming $\int\,dx$, then power to $\frac{1}{p}$, and finally summing $\int \, dy$. Minkowski inequality is especially useful for functions of the form 
	\[f(x, y) = \phi_\epsilon (x - y) \, u(y)\]
	where $\phi_\epsilon$ is some sort of bump function $(\phi_\epsilon (x) \rightarrow \delta (x)$ as $\epsilon \rightarrow 0$. Then, defining the convolution operator 
	\[(J_\epsilon u) (x) = \int \varphi_\epsilon (x - y) u(y)\, dy\]
	Now we can approximate the delta function $\delta (x - y)$ with $C^\infty$ bump functions by "squeezing them" as $\phi_\epsilon (t) = \phi(t / \epsilon) \epsilon^{-n}$ (we power to $n$ to normalize over $t \in \mathbb{R}^n$, an $n$-dimensional input). Therefore, we can approximate $u$ with arbitrary accuracy with a $C^\infty$ function. 
	\[\underbrace{(J_\epsilon u) (x) = \int \varphi_\epsilon (x - y) u(y)\, dy}_{\text{smooth}} \approx \int \delta (x - y) u(y) \,dy = u(x)\]
	Assuming that some smooth and compactly supported $\varphi \in C^\infty_C (\mathbb{R}^n)$ exists (we can construct this), we can construct a family of approximate identities/deltas (since deltas are identities for convolutions), called \textit{mollifiers} (which "mollify" or smooth out an arbitrary function) 
	\[\{\varphi_\epsilon\}_\epsilon \text{ s.t. } \varphi_\epsilon (t) = \epsilon^{-n} \varphi(t / \epsilon)\]
	We have 
	\[(J_\epsilon u) (x) \coloneqq (\varphi_\epsilon \ast u) (x) = \int \varphi_\epsilon (x - y) u(y) \, dy \rightarrow u(x) \text{ in } L^p\]
	So, $||(J_\epsilon u (x) - u(x)||_{L^p} \rightarrow 0$ using Minkowski inequality. This property is very useful because if we're getting an arbitrary signal, we can interpret it as a function, assume the white noise is Gaussian, and then the convolution would be a smooth signal. So in reality, we don't have to worry about smoothness. 

	\begin{theorem}
	If there exists some function $u$ such that 
	\[\int u \varphi = 0 \text{ for all } \phi \in C^\infty_C\]
	Then $u = 0$ almost everywhere. 
	\end{theorem}

